% STATUS: COMPLETE DRAFT
% Parte 1: numeri, frazioni e potenze, equazioni, proporzioni e percentuali, statistica e probabilità
% Every number is checked in verifica.py (same ids).

% ================================================================
\sezione{Numeri naturali e proprietà}{Natural numbers and properties}
% ================================================================

\begin{esercizio}{N1}{1}{Le proprietà delle operazioni}{Properties of the operations}
Scrivi il nome della proprietà usata in ogni uguaglianza.
\en{Write the name of the property used in each equality.}
\begin{multicols}{2}
\begin{enumerate}
\item $7\cdot(10+3)=7\cdot10+7\cdot3$
\item $36:4=(36:2):(4:2)$
\item $5\cdot 0\cdot 8=0$
\item $12+8=8+12$
\item $(4\cdot 5)\cdot 2=4\cdot(5\cdot 2)$
\end{enumerate}
\end{multicols}
\vspace{-4pt}
Poi calcola, se è possibile, $0:9$, \ $9:0$, \ $0:0$ e spiega la risposta.
\en{Then compute, if possible, $0:9$, $9:0$, $0:0$ and explain your answer.}
\end{esercizio}
\begin{solbox}
\begin{enumerate}
\item Proprietà \textbf{distributiva} della moltiplicazione rispetto all'addizione. \enx{distributive}
\item Proprietà \textbf{invariantiva} della divisione: dividendo e divisore divisi per lo stesso numero. \enx{invariance of the quotient}
\item Legge di \textbf{annullamento del prodotto}: se un fattore è $0$, il prodotto è $0$. \enx{zero product law}
\item Proprietà \textbf{commutativa} dell'addizione. \enx{commutative}
\item Proprietà \textbf{associativa} della moltiplicazione. \enx{associative}
\end{enumerate}
$0:9=0$, perché $0\cdot 9=0$.\quad
$9:0$ è \textbf{impossibile}: nessun numero moltiplicato per $0$ dà $9$.\quad
$0:0$ è \textbf{indeterminata}: ogni numero moltiplicato per $0$ dà $0$.
\en{$0:9=0$. $9:0$ is impossible (no number times $0$ gives $9$). $0:0$ is undetermined (every number works).}
\end{solbox}

\begin{esercizio}{N2}{2}{mcm e MCD nei problemi}{LCM and GCD in word problems}
\begin{enumerate}
\item Due autobus partono insieme dal capolinea alle 7:00. Il primo ripassa ogni 12 minuti, il secondo ogni 18 minuti. A che ora ripartono di nuovo insieme per la prima volta?
\en{Two buses leave the terminal together at 7:00. One comes back every 12 minutes, the other every 18 minutes. When do they leave together again for the first time?}
\item Una maestra ha 48 penne e 72 matite. Vuole preparare astucci tutti uguali, usando tutto il materiale. Qual è il numero massimo di astucci? Quante penne e quante matite ci sono in ognuno?
\en{A teacher has 48 pens and 72 pencils. She wants identical pencil cases, using everything. What is the largest number of cases? How many pens and pencils in each?}
\end{enumerate}
Scomponi prima i numeri in fattori primi. \enx{First write each number as a product of primes.}
\end{esercizio}
\begin{solbox}
$12=2^2\cdot3$, \quad $18=2\cdot3^2$, \quad $48=2^4\cdot3$, \quad $72=2^3\cdot3^2$.
\begin{enumerate}
\item Cerco quando gli eventi si ripetono \emph{insieme}: $\text{mcm}(12,18)=2^2\cdot3^2=36$. Ripartono insieme dopo 36 minuti, alle \textbf{7:36}.
\en{Events repeating together: least common multiple. They meet after 36 minutes, at 7:36.}
\item Divido in parti uguali \emph{il più grandi possibile}: $\text{MCD}(48,72)=2^3\cdot3=24$. Quindi \textbf{24 astucci}, ognuno con $48:24=2$ penne e $72:24=3$ matite.
\en{Sharing into the largest equal groups: greatest common divisor, 24 cases with 2 pens and 3 pencils each.}
\end{enumerate}
\errore{mcm: fattori comuni e non comuni con l'esponente \emph{maggiore}; MCD: solo i fattori comuni con l'esponente \emph{minore}.}
{LCM takes all prime factors with the highest power; GCD takes only common factors with the lowest power.}
\end{solbox}

% ================================================================
\sezione{Frazioni, potenze e radici}{Fractions, powers and roots}
% ================================================================

\begin{esercizio}{F1}{1}{Espressione con le frazioni}{Expression with fractions}
Calcola: \enx{Compute:}
\[\left(\frac34+\frac16\right)\cdot\frac{6}{11}-\frac13\]
\end{esercizio}
\begin{solbox}
\[
\left(\frac{9}{12}+\frac{2}{12}\right)\cdot\frac{6}{11}-\frac13
=\frac{\cancel{11}}{\cancelto{2}{12}}\cdot\frac{\cancel{6}}{\cancel{11}}-\frac13
=\frac12-\frac13=\frac{3-2}{6}=\boxed{\frac16}
\]
Prima le parentesi, poi la moltiplicazione, infine la sottrazione.
\en{First the brackets, then the multiplication, and the subtraction last.}
\end{solbox}

\begin{esercizio}{F2}{1}{Le proprietà delle potenze}{Properties of powers}
Calcola usando le proprietà delle potenze. \enx{Compute using the properties of powers.}
\begin{multicols}{2}
\begin{enumerate}
\item $2^3\cdot 2^4:2^5$
\item $(3^2)^3:3^4$
\item $\left(\dfrac25\right)^3\cdot\left(\dfrac52\right)^3$
\item $(-2)^3,\ \ (-2)^4,\ \ -2^4,\ \ \left(-\dfrac13\right)^2$
\item $7^0,\ \ 1^9,\ \ 0^5$
\end{enumerate}
\end{multicols}
\end{esercizio}
\begin{solbox}
\begin{enumerate}
\item Stessa base: si sommano e sottraggono gli esponenti. $2^{3+4-5}=2^2=\mathbf{4}$
\item Potenza di potenza: si moltiplicano gli esponenti. $3^{6}:3^4=3^{2}=\mathbf{9}$
\item Stesso esponente: si moltiplicano le basi. $\left(\frac25\cdot\frac52\right)^3=1^3=\mathbf{1}$
\item $(-2)^3=\mathbf{-8}$ (esponente dispari: il segno resta); $(-2)^4=\mathbf{16}$ (esponente pari: positivo); $-2^4=\mathbf{-16}$ (il meno è \emph{fuori} dalla potenza); $\left(-\frac13\right)^2=\mathbf{\frac19}$.
\item $7^0=\mathbf{1}$ (ogni numero diverso da zero elevato a $0$ fa $1$); $1^9=\mathbf{1}$; $0^5=\mathbf{0}$.
\end{enumerate}
\en{Same base: add or subtract exponents. Power of a power: multiply exponents. Same exponent: multiply the bases. An odd exponent keeps the minus sign; an even one makes the result positive; in $-2^4$ the minus sign is outside the power.}
\end{solbox}

\begin{esercizio}{F3}{2}{Potenze di frazioni}{Powers of fractions}
Calcola: \enx{Compute:}
\[\left[\left(\frac53\right)^4:\left(\frac53\right)^2\right]\cdot\frac{9}{25}+\left(\frac12\right)^2\]
\end{esercizio}
\begin{solbox}
\[
\left(\frac53\right)^{4-2}\cdot\frac{9}{25}+\frac14
=\frac{25}{9}\cdot\frac{9}{25}+\frac14
=1+\frac14=\boxed{\frac54}
\]
\end{solbox}

\begin{esercizio}{F4}{2}{Numeri decimali periodici e radici quadrate}{Repeating decimals and square roots}
\begin{enumerate}
\item Trova la frazione generatrice di $0,\overline{3}$, \ $0,\overline{45}$, \ $1,2\overline{5}$.
\en{Write each repeating decimal as a fraction.}
\item Calcola $\sqrt{144}$, \ $\sqrt{\dfrac{9}{25}}$, \ $\sqrt{0,49}$.
\en{Compute the square roots.}
\item Tra quali due numeri interi consecutivi si trova $\sqrt{50}$?
\en{Between which two consecutive whole numbers is $\sqrt{50}$?}
\end{enumerate}
\end{esercizio}
\begin{solbox}
\begin{enumerate}
\item Regola: al numeratore tutto il numero senza virgola meno la parte che non si ripete; al denominatore tanti $9$ quante sono le cifre del periodo, seguiti da tanti $0$ quante sono le cifre dell'antiperiodo.
\en{Numerator: the whole number without the comma minus the non repeating part. Denominator: one 9 for each repeating digit, then one 0 for each non repeating decimal digit.}
\[0,\overline{3}=\frac39=\frac13\qquad 0,\overline{45}=\frac{45}{99}=\frac{5}{11}\qquad 1,2\overline{5}=\frac{125-12}{90}=\frac{113}{90}\]
\item $\sqrt{144}=\mathbf{12}$ perché $12^2=144$; \quad $\sqrt{\frac{9}{25}}=\mathbf{\frac35}$; \quad $\sqrt{0,49}=\mathbf{0,7}$ perché $0,7^2=0,49$.
\item $7^2=49<50<64=8^2$, quindi $\sqrt{50}$ è tra \textbf{7 e 8} (circa $7,07$).
\end{enumerate}
\end{solbox}

\begin{esercizio}{F5}{3}{Espressione lunga}{A long expression}
Calcola con attenzione ai segni. \enx{Compute, paying attention to the signs.}
\[\left\{\left[\left(1-\frac14\right)\cdot\left(1+\frac13\right)-\left(\frac12\right)^2\right]:\left(\frac34\right)^2+\left(-\frac23\right)^2\right\}\cdot\frac{9}{16}-\left(-\frac12\right)^3\]
\end{esercizio}
\begin{solbox}
Tonde, poi quadre, poi graffe. \enx{Round brackets, then square, then curly.}
\begin{align*}
&\left(1-\tfrac14\right)\cdot\left(1+\tfrac13\right)=\tfrac34\cdot\tfrac43=1
&&\left[\,1-\tfrac14\,\right]=\tfrac34\\
&\tfrac34:\left(\tfrac34\right)^2=\tfrac34\cdot\tfrac{16}{9}=\tfrac43
&&\tfrac43+\left(-\tfrac23\right)^2=\tfrac43+\tfrac49=\tfrac{12+4}{9}=\tfrac{16}{9}\\
&\tfrac{16}{9}\cdot\tfrac{9}{16}=1
&&1-\left(-\tfrac12\right)^3=1-\left(-\tfrac18\right)=1+\tfrac18=\boxed{\tfrac98}
\end{align*}
\errore{$\left(-\frac12\right)^3=-\frac18$, e \emph{meno per meno} fa più: $1-\left(-\frac18\right)=1+\frac18$.}
{a negative base with an odd exponent stays negative, and subtracting a negative number means adding.}
\end{solbox}

% ================================================================
\sezione{Equazioni di primo grado}{Linear equations}
% ================================================================

\begin{esercizio}{E1}{1}{Equazioni semplici}{Simple equations}
Risolvi e fai la verifica. \enx{Solve and check.}
\begin{multicols}{2}
\begin{enumerate}
\item $3x-7=2x+5$
\item $4(x-2)=2x+6$
\end{enumerate}
\end{multicols}
\end{esercizio}
\begin{solbox}
\begin{enumerate}
\item $3x-2x=5+7\ \Rightarrow\ \boxed{x=12}$. \ Verifica: $3\cdot12-7=29$ e $2\cdot12+5=29$. \checkmark
\item $4x-8=2x+6\ \Rightarrow\ 2x=14\ \Rightarrow\ \boxed{x=7}$. \ Verifica: $4\cdot5=20$ e $14+6=20$. \checkmark
\end{enumerate}
\en{Move the $x$ terms to one side and the numbers to the other, changing sign when a term crosses the equals sign.}
\end{solbox}

\begin{esercizio}{E2}{2}{Equazione con le parentesi}{Equation with brackets}
Risolvi e verifica: \enx{Solve and check:}
\[2(x+3)-3(x-1)=5-2x\]
\end{esercizio}
\begin{solbox}
\[2x+6-3x+3=5-2x\ \Rightarrow\ -x+9=5-2x\ \Rightarrow\ -x+2x=5-9\ \Rightarrow\ \boxed{x=-4}\]
Verifica: $2(-1)-3(-5)=-2+15=13$ \ e \ $5-2(-4)=5+8=13$. \checkmark
\errore{$-3(x-1)=-3x\mathbf{+3}$: il meno davanti alla parentesi cambia \emph{tutti} i segni dentro.}
{the minus sign in front of the bracket changes every sign inside it.}
\end{solbox}

\begin{esercizio}{E3}{2}{Equazione con le frazioni}{Equation with fractions}
Risolvi e verifica: \enx{Solve and check:}
\[\frac{x-1}{2}+\frac{x}{3}=2+\frac{x+1}{6}\]
\end{esercizio}
\begin{solbox}
Moltiplico tutto per il mcm dei denominatori, $6$: \enx{multiply every term by 6}
\[3(x-1)+2x=12+(x+1)\ \Rightarrow\ 3x-3+2x=x+13\ \Rightarrow\ 4x=16\ \Rightarrow\ \boxed{x=4}\]
Verifica: $\frac32+\frac43=\frac{9+8}{6}=\frac{17}{6}$ \ e \ $2+\frac56=\frac{17}{6}$. \checkmark
\errore{anche il $2$ va moltiplicato per $6$.}{the whole number $2$ must be multiplied by 6 too.}
\end{solbox}

\begin{esercizio}{E4}{3}{Problemi risolti con le equazioni}{Word problems with equations}
\begin{enumerate}
\item Il perimetro di un rettangolo è $56\cm$ e la base supera l'altezza di $8\cm$. Trova base, altezza e area.
\en{A rectangle has perimeter $56\cm$ and its base is $8\cm$ longer than its height. Find base, height and area.}
\item La somma di tre numeri naturali consecutivi è $84$. Quali sono?
\en{The sum of three consecutive natural numbers is $84$. Find them.}
\item Una di queste equazioni è impossibile e l'altra è indeterminata. Quale è quale? \ $2(x+1)=2x+2$ \quad $3x+1=3x+5$
\en{One equation is impossible and the other is undetermined. Which is which?}
\end{enumerate}
\end{esercizio}
\begin{solbox}
\begin{enumerate}
\item Altezza $x$, base $x+8$: \ $2(x+x+8)=56\Rightarrow 4x+16=56\Rightarrow x=10$. \ Altezza $\mathbf{10\cm}$, base $\mathbf{18\cm}$, area $10\cdot18=\mathbf{180\cmq}$.
\item $n+(n+1)+(n+2)=84\Rightarrow 3n+3=84\Rightarrow n=27$. \ I numeri sono $\mathbf{27,\ 28,\ 29}$.
\item $2x+2=2x+2\Rightarrow 0x=0$: vera per ogni $x$, \textbf{indeterminata}. \quad $3x+1=3x+5\Rightarrow 0x=4$: nessun $x$, \textbf{impossibile}.
\en{$0x=0$ holds for every $x$ (undetermined); $0x=4$ never holds (impossible).}
\end{enumerate}
\end{solbox}

% ================================================================
\sezione{Proporzioni e percentuali}{Proportions and percentages}
% ================================================================

\begin{esercizio}{P1}{1}{Proporzioni}{Proportions}
Trova il termine incognito. \enx{Find the unknown term.}
\begin{multicols}{3}
\begin{enumerate}
\item $6:x=9:12$
\item $x:5=14:10$
\item $(x+2):x=5:3$
\end{enumerate}
\end{multicols}
\end{esercizio}
\begin{solbox}
Proprietà fondamentale: prodotto dei medi $=$ prodotto degli estremi. \enx{product of the means equals product of the extremes}
\begin{enumerate}
\item $9x=6\cdot12=72\Rightarrow \boxed{x=8}$
\item $10x=5\cdot14=70\Rightarrow \boxed{x=7}$
\item $3(x+2)=5x\Rightarrow 3x+6=5x\Rightarrow 6=2x\Rightarrow \boxed{x=3}$. \ Verifica: $5:3=5:3$. \checkmark
\end{enumerate}
\end{solbox}

\begin{esercizio}{P2}{1}{Calcolo di percentuali}{Working with percentages}
\begin{enumerate}
\item Quanto è il $15\%$ di $240$? \enx{What is 15\% of 240?}
\item $18$ è quale percentuale di $72$? \enx{18 is what percentage of 72?}
\item Il $30\%$ di un numero è $45$. Qual è il numero? \enx{30\% of a number is 45. Find the number.}
\end{enumerate}
\end{esercizio}
\begin{solbox}
\begin{enumerate}
\item $240\cdot\frac{15}{100}=\mathbf{36}$
\item $\frac{18}{72}\cdot100=\frac14\cdot100=\mathbf{25\%}$
\item $45:30=x:100\Rightarrow x=\frac{45\cdot100}{30}=\mathbf{150}$
\end{enumerate}
\end{solbox}

\begin{esercizio}{P3}{2}{Sconti e aumenti}{Discounts and increases}
\begin{enumerate}
\item Un paio di scarpe costa $170$\,€ dopo uno sconto del $20\%$. Qual era il prezzo prima dello sconto?
\en{A pair of shoes costs 170\,€ after a 20\% discount. What was the price before the discount?}
\item Un giubbotto costa $80$\,€. Il prezzo aumenta del $10\%$ e poi viene scontato del $10\%$. Torna a costare $80$\,€? Perché?
\en{A jacket costs 80\,€. Its price goes up by 10\%, then goes down by 10\%. Is it 80\,€ again? Why?}
\end{enumerate}
\end{esercizio}
\begin{solbox}
\begin{enumerate}
\item Dopo lo sconto si paga $100\%-20\%=80\%$ del prezzo iniziale:
\[80:170=100:x\ \Rightarrow\ x=\frac{170\cdot100}{80}=\boxed{212,50\ \text{€}}\qquad\text{verifica: } 212,50\cdot0,8=170\ \checkmark\]
\errore{$170+20\%$ di $170=204$ è sbagliato: il $20\%$ si calcola sul prezzo \emph{iniziale}, non su quello scontato.}
{adding 20\% of 170 is wrong, because the 20\% was taken from the original price.}
\item $80\cdot1,1=88$\,€, \ poi $88\cdot0,9=\mathbf{79,20}$\,€. \ \textbf{No}: il secondo $10\%$ si calcola su $88$, un numero più grande, quindi toglie più di quanto era stato aggiunto.
\en{No: the second 10\% is taken from 88, so it removes more than was added.}
\end{enumerate}
\end{solbox}

\begin{esercizio}{P4}{2}{Consumo di energia}{Energy use}
Un elettrodomestico consuma $1,5$\,kWh in $4$ ore.
\en{An appliance uses 1.5\,kWh in 4 hours.}
\begin{enumerate}
\item Quanto consuma in $10$ ore? \enx{How much in 10 hours?}
\item E se lavora per $10$ ore al $60\%$ della potenza? \enx{And at 60\% of its power, for 10 hours?}
\item Se $1$\,kWh costa $0,28$\,€, quanto si spende nel caso b)? \enx{If 1\,kWh costs 0.28\,€, what does case b) cost?}
\end{enumerate}
\end{esercizio}
\begin{solbox}
\begin{enumerate}
\item Proporzionalità diretta: $1,5:4=x:10\Rightarrow x=\frac{1,5\cdot10}{4}=\mathbf{3,75}$\,kWh.
\item $3,75\cdot\frac{60}{100}=\mathbf{2,25}$\,kWh.
\item $2,25\cdot0,28=\mathbf{0,63}$\,€.
\end{enumerate}
\end{solbox}

\begin{esercizio}{P5}{3}{Proporzionalità diretta e inversa}{Direct and inverse proportion}
\begin{enumerate}
\item Un atleta percorre $12$\,km in $48$ minuti, a velocità costante. In quanto tempo percorre $20$\,km?
\en{A runner covers 12\,km in 48 minutes at constant speed. How long for 20\,km?}
\item $6$ operai costruiscono un muro in $10$ giorni. Quanti giorni servono a $4$ operai che lavorano allo stesso ritmo?
\en{6 workers build a wall in 10 days. How many days for 4 workers at the same pace?}
\item Per ogni caso, le grandezze sono direttamente o inversamente proporzionali? Scrivi la costante.
\en{In each case, are the quantities directly or inversely proportional? Write the constant.}
\end{enumerate}
\end{esercizio}
\begin{solbox}
\begin{enumerate}
\item $12:48=20:x\Rightarrow x=\frac{48\cdot20}{12}=\mathbf{80}$ minuti $=1$\,h $20$\,min.
\item Meno operai, \emph{più} giorni: $6\cdot10=4\cdot x\Rightarrow x=\mathbf{15}$ giorni.
\item Caso a) \textbf{diretta}: il rapporto è costante, $\frac{48}{12}=4$ minuti per km. \ Caso b) \textbf{inversa}: il prodotto è costante, $6\cdot10=60$.
\en{a) direct proportion, constant ratio 4 minutes per km; b) inverse proportion, constant product 60.}
\end{enumerate}
\errore{nel caso b) la proporzione $6:10=4:x$ è sbagliata: darebbe meno giorni con meno operai.}
{in case b) the direct proportion $6:10=4:x$ is wrong: fewer workers would finish sooner.}
\end{solbox}

% ================================================================
\sezione{Statistica e probabilità}{Statistics and probability}
% ================================================================

\begin{esercizio}{S1}{1}{Media, moda, mediana}{Mean, mode, median}
I voti di Luca in matematica sono: \ $6,\ 7,\ 5,\ 8,\ 7,\ 9,\ 7,\ 6$. \ Calcola media, moda, mediana e campo di variazione.
\en{Luca's maths marks are listed above. Find the mean, mode, median and range.}
\end{esercizio}
\begin{solbox}
Media: $\frac{6+7+5+8+7+9+7+6}{8}=\frac{55}{8}=\mathbf{6,875}$.\quad
Dati in ordine: $5,6,6,7,7,7,8,9$: \ mediana $=\frac{7+7}{2}=\mathbf{7}$ (i due valori centrali).\quad
Moda $=\mathbf{7}$ (compare 3 volte).\quad
Campo di variazione $=9-5=\mathbf{4}$.
\en{With an even number of data the median is the mean of the two middle values, after sorting.}
\end{solbox}

\begin{esercizio}{S2}{2}{Dalla tabella all'areogramma}{From a table to a pie chart}
In una biblioteca ci sono: gialli $96$, fantasy $60$, romanzi storici $48$, biografie $36$.
\en{A library has 96 crime novels, 60 fantasy, 48 historical novels and 36 biographies.}
\begin{enumerate}
\item Calcola il totale e la frequenza relativa percentuale di ogni genere.
\en{Find the total and the percentage of each type.}
\item Calcola l'ampiezza in gradi di ogni settore e disegna l'areogramma con il goniometro.
\en{Find the angle of each sector and draw the pie chart with a protractor.}
\end{enumerate}
\end{esercizio}
\begin{solbox}
Totale $96+60+48+36=240$.\quad Percentuale $=\frac{\text{frequenza}}{240}\cdot100$; \quad gradi $=\frac{\text{frequenza}}{240}\cdot360^\circ$.
\begin{center}
\begin{tabular}{@{}lccc@{}}
\toprule
genere & frequenza & percentuale & gradi\\
\midrule
gialli & $96$ & $40\%$ & $144^\circ$\\
fantasy & $60$ & $25\%$ & $90^\circ$\\
romanzi storici & $48$ & $20\%$ & $72^\circ$\\
biografie & $36$ & $15\%$ & $54^\circ$\\
\midrule
totale & $240$ & $100\%$ & $360^\circ$\\
\bottomrule
\end{tabular}
\hspace{1cm}
% sectors from 90 deg clockwise: 144, 90, 72, 54 -> boundaries 90, -54, -144, -216, -270
\begin{tikzpicture}[baseline=(c.base)]
  \coordinate (c) at (0,0);
  \def\R{1.35}
  \fill[fgA!55] (0,0) -- (90:\R) arc (90:-54:\R) -- cycle;
  \fill[fgB!50] (0,0) -- (-54:\R) arc (-54:-144:\R) -- cycle;
  \fill[fgC!50] (0,0) -- (-144:\R) arc (-144:-216:\R) -- cycle;
  \fill[fgD!55] (0,0) -- (-216:\R) arc (-216:-270:\R) -- cycle;
  \foreach \a in {90,-54,-144,-216} \draw[white, line width=1pt] (0,0) -- (\a:\R);
  \node[font=\scriptsize] at (18:0.85) {$144^\circ$};
  \node[font=\scriptsize] at (-99:0.85) {$90^\circ$};
  \node[font=\scriptsize] at (-180:0.85) {$72^\circ$};
  \node[font=\scriptsize] at (-243:0.9) {$54^\circ$};
\end{tikzpicture}
\end{center}
Controllo: la somma delle percentuali è $100\%$ e quella dei gradi è $360^\circ$. \enx{check both totals}
\end{solbox}

\begin{esercizio}{S3}{2}{Dall'areogramma ai dati}{From a pie chart back to the data}
L'areogramma rappresenta i giudizi di $72$ alunni. \enx{The pie chart shows the grades of 72 pupils.}
\begin{center}
% sectors: sufficiente 150, buono 120, ottimo (unknown) 90 ; drawn to scale
\begin{tikzpicture}
  \def\R{1.5}
  \fill[fgB!40] (0,0) -- (90:\R) arc (90:-60:\R) -- cycle;
  \fill[fgC!40] (0,0) -- (-60:\R) arc (-60:-180:\R) -- cycle;
  \fill[fgD!45] (0,0) -- (-180:\R) arc (-180:-270:\R) -- cycle;
  \foreach \a in {90,-60,-180} \draw[white, line width=1pt] (0,0) -- (\a:\R);
  \node[font=\small, align=center] at (15:0.85) {sufficiente\\$150^\circ$};
  \node[font=\small, align=center] at (-120:0.85) {buono\\$120^\circ$};
  \node[font=\small, align=center] at (-225:0.8) {ottimo\\$?$};
\end{tikzpicture}
\end{center}
\begin{enumerate}
\item Quanto è ampio il settore «ottimo»? \enx{What is the angle of the sector ``excellent''?}
\item Quanti alunni hanno preso ciascun giudizio? \enx{How many pupils got each grade?}
\item Che percentuale di alunni ha preso «ottimo»? \enx{What percentage got ``excellent''?}
\end{enumerate}
\end{esercizio}
\begin{solbox}
\begin{enumerate}
\item $360^\circ-(150^\circ+120^\circ)=\mathbf{90^\circ}$.
\item Alunni $=\frac{\text{gradi}}{360^\circ}\cdot72$: \ sufficiente $\frac{150}{360}\cdot72=\mathbf{30}$; \ buono $\frac{120}{360}\cdot72=\mathbf{24}$; \ ottimo $\frac{90}{360}\cdot72=\mathbf{18}$. \ Controllo: $30+24+18=72$. \checkmark
\item $\frac{90}{360}=\frac14=\mathbf{25\%}$.
\end{enumerate}
\end{solbox}

\begin{esercizio}{S4}{1}{Probabilità con un dado}{Probability with one die}
Lanci un dado a sei facce. Scrivi la probabilità come frazione, numero decimale e percentuale.
\en{You roll a six sided die. Write each probability as a fraction, a decimal and a percentage.}
\begin{multicols}{2}
\begin{enumerate}
\item Esce un numero pari. \enx{an even number}
\item Esce un numero maggiore di $4$. \enx{more than 4}
\item Esce il $7$. \enx{a 7}
\item Esce un numero da $1$ a $6$. \enx{a number from 1 to 6}
\end{enumerate}
\end{multicols}
\end{esercizio}
\begin{solbox}
$P=\dfrac{\text{casi favorevoli}}{\text{casi possibili}}$
\begin{enumerate}
\item $\{2,4,6\}$: \ $\frac36=\frac12=0,5=\mathbf{50\%}$
\item $\{5,6\}$: \ $\frac26=\frac13\approx0,33\approx\mathbf{33,3\%}$
\item Evento \textbf{impossibile}: $P=\mathbf{0}$.
\item Evento \textbf{certo}: $P=\mathbf{1}=100\%$.
\end{enumerate}
\en{Probability is always between 0 (impossible) and 1 (certain).}
\end{solbox}

\begin{esercizio}{S5}{3}{Urne e due dadi}{Urns and two dice}
\begin{enumerate}
\item Un'urna contiene $5$ palline rosse, $3$ blu e $2$ verdi. Ne estrai una. Calcola la probabilità che sia: rossa; non blu; rossa oppure verde.
\en{An urn holds 5 red, 3 blue and 2 green balls. You draw one. Find P(red), P(not blue), P(red or green).}
\item Lanci due dadi. Costruisci la tabella delle somme e calcola la probabilità che: la somma sia $7$; i due numeri siano uguali; la somma sia maggiore di $10$.
\en{You roll two dice. Build the table of sums and find P(sum 7), P(both the same), P(sum greater than 10).}
\end{enumerate}
\end{esercizio}
\begin{solbox}
\begin{enumerate}
\item Casi possibili $10$. \ $P(\text{rossa})=\frac{5}{10}=\mathbf{\frac12}$; \ $P(\text{non blu})=1-\frac{3}{10}=\mathbf{\frac{7}{10}}$ (evento contrario); \ $P(\text{rossa o verde})=\frac{5+2}{10}=\mathbf{\frac{7}{10}}$.
\item $6\cdot6=36$ casi possibili.
\begin{center}
% 6x6 table of sums; sum 7 in green, doubles in orange (4+... none overlap), sum > 10 in red
\begin{tikzpicture}[scale=0.52]
  \foreach \i in {1,...,6} {
    \node[font=\scriptsize\bfseries] at (\i,0.1) {\i};
    \node[font=\scriptsize\bfseries] at (-0.05,-\i) {\i};
    \foreach \j in {1,...,6} {
      \pgfmathtruncatemacro{\s}{\i+\j}
      \ifnum\s=7 \fill[fgC!35] (\i-0.5,-\j-0.5) rectangle (\i+0.5,-\j+0.5); \fi
      \ifnum\i=\j \fill[fgD!40] (\i-0.5,-\j-0.5) rectangle (\i+0.5,-\j+0.5); \fi
      \ifnum\s>10 \draw[fgA, line width=1pt] (\i-0.42,-\j-0.42) rectangle (\i+0.42,-\j+0.42); \fi
      \node[font=\scriptsize] at (\i,-\j) {\s};
    }
  }
  \foreach \k in {0,...,6} {\draw[black!40] (\k+0.5,-0.5) -- (\k+0.5,-6.5); \draw[black!40] (0.5,-\k-0.5) -- (6.5,-\k-0.5);}
  \node[font=\scriptsize, anchor=west] at (7,-1.5) {\textcolor{fgC!80!black}{$\blacksquare$} somma 7: $6$ casi};
  \node[font=\scriptsize, anchor=west] at (7,-2.7) {\textcolor{fgD}{$\blacksquare$} numeri uguali: $6$ casi};
  \node[font=\scriptsize, anchor=west] at (7,-3.9) {\textcolor{fgA}{$\square$} somma $>10$: $3$ casi};
\end{tikzpicture}
\end{center}
$P(\text{somma }7)=\frac{6}{36}=\mathbf{\frac16}$; \quad $P(\text{uguali})=\frac{6}{36}=\mathbf{\frac16}$; \quad $P(\text{somma}>10)=\frac{3}{36}=\mathbf{\frac{1}{12}}$ (casi $5{+}6$, $6{+}5$, $6{+}6$).
\end{enumerate}
\errore{$1+6$ e $6+1$ sono due casi diversi: nella tabella occupano due caselle.}{$1+6$ and $6+1$ are two different outcomes.}
\end{solbox}
